\documentclass[11pt]{article}
\usepackage[margin=1in]{geometry}
\usepackage[T1]{fontenc}
\usepackage{amsmath,amssymb,amsthm,mathtools}
\usepackage[hidelinks]{hyperref}
\newtheorem{theorem}{Theorem}
\newtheorem{lemma}[theorem]{Lemma}
\newtheorem{proposition}[theorem]{Proposition}
\newtheorem{corollary}[theorem]{Corollary}
\theoremstyle{definition}
\newtheorem{definition}[theorem]{Definition}
\newtheorem{remark}[theorem]{Remark}
\newcommand{\Th}{\text{\th}}
\title{Mode stability of subextremal Kerr--de Sitter black holes\\[2pt]
\large Teukolsky ($s=\pm2$) and ungauged linearised Einstein equations: the proof as it stands}
\author{Jaideep Sai Padhi}
\date{8 November 2026 --- working version}
\begin{document}
\maketitle

\begin{abstract}
For every subextremal Kerr--de Sitter (KdS) black hole we prove (A) that the spin $\pm2$ Teukolsky equation has no
nonzero future-regular mode solutions with $\operatorname{Im}\sigma\ge0$, apart from the allowed stationary ones, and
(B) that every generalised mode solution of the ungauged linearised Einstein equation $D\mathrm{Ric}(h)=\Lambda h$ with
$\operatorname{Im}\sigma\ge0$ is pure gauge if $\sigma\neq0$, and a linearised KdS metric plus pure gauge if $\sigma=0$.
Statement (B) is the hypothesis of the full-subextremal nonlinear stability theorem announced by Hintz--Petersen--Vasy.
\textbf{Status.} Both parts are proved. Part (B) is proved given (A); in the large-$A$ regime R3 of the event lemma
the multiplier is constructed by saturating one condition and solving the other with equality
(Section~\ref{sec:open}), following the structure of Hintz's Klein--Gordon argument. Every step marked \emph{[cert]} uses an exact symbolic or a short
numerical verification; the scripts are named in the repository ledger \texttt{kds/LEMMAS.md} (item numbers below).
\end{abstract}

\tableofcontents

\section{Setting and statements}
KdS with parameters $(\Lambda,M,a)$, $\Lambda>0$, in Boyer--Lindquist coordinates has
$\Delta_r=(r^2+a^2)(1-\Lambda r^2/3)-2Mr=-\tfrac\Lambda3\prod_{j=0}^3(r-r_j)$, $\alpha=\Lambda a^2/3$, $\Xi=1+\alpha$.
\emph{Subextremal} means four distinct real roots $r_3<0<r_0<r_1=r_e<r_2=r_c$; then $\sum r_j=0$.
Modes are $e^{-i\sigma t+im\varphi}$ times functions of $(r,\theta)$, smooth across both horizons in the regular time $t_*$.

\begin{theorem}[A]\label{thmA}
For every subextremal KdS the spin $s=\pm2$ Teukolsky equation has no nonzero future-regular mode with
$\operatorname{Im}\sigma\ge0$, except those with $\sigma=0$ permitted below.
\end{theorem}

\begin{theorem}[B]\label{thmB}
Let $h$ be a generalised mode, $(L_T+i\sigma)^k h=0$ with $\operatorname{Im}\sigma\ge0$, of $D\mathrm{Ric}(h)=\Lambda h$,
smooth across both horizons. If $\sigma\neq0$ then $h=\mathcal L_Xg$ with $X$ smooth; if $\sigma=0$ then
$h=c_M\partial_Mg+c_a\partial_ag+\mathcal L_Xg$ with $X$ polynomial in $t_*$.
\end{theorem}

\section{Part A: Teukolsky mode stability}
\subsection{Heun form and the window}
The M\"obius map $r_0,r_1,r_2,r_3\mapsto0,1,t,\infty$ puts the radial equation in Heun form with
$\mu(z)=z(z-1)(t-z)$, $t=\frac{(r_1-r_3)(r_2-r_0)}{(r_1-r_0)(r_2-r_3)}>1$ (item 26). With
$K(r)=\omega(r^2+a^2)-am$ and $\kappa_j=i\Xi K(r_j)/\Delta_r'(r_j)$, put $e_1=2(\kappa_1+\kappa_3)$, $e_t=2(\kappa_2+\kappa_3)$.
For real $\omega$ the dual equation is real and the Wronskian identity gives
$\operatorname{Im}e_1\,|c_1|^2=t\operatorname{Im}e_t\,|c_t|^2$ (item 34).

\begin{proposition}[the window; item 88]\label{prop:window}
Let $m\neq0$, $w=\omega/m$, $\rho_j=r_j^2+a^2$, $\Omega_j=a/\rho_j$, $D_j=|\Delta_r'(r_j)|$, $P_j=\rho_j/D_j$. Then
$\operatorname{Im}e_1=2\Xi m f_1(w)$, $\operatorname{Im}e_t=2\Xi m f_2(w)$ with
$f_1=P_1(w-\Omega_1)+P_3(w-\Omega_3)$ increasing and $f_2=-P_2(w-\Omega_2)+P_3(w-\Omega_3)$ decreasing; their roots
satisfy $0<w_2<w_1$. Hence $\operatorname{Im}e_1\cdot\operatorname{Im}e_t>0$ exactly on the open window $(w_2,w_1)$.
\end{proposition}
\begin{proof}
$\operatorname{sign}\Delta_r'(r_j)=(-,+,-,+)$. Since $\rho/\Delta_r$ and $1/\Delta_r$ have zero residue at infinity,
$P_1+P_3=P_0+P_2$ and $D_1^{-1}+D_3^{-1}=D_0^{-1}+D_2^{-1}$. So the slope of $f_2$ is $P_0-P_1$, and $P_1>P_0$ because
$D_0/D_1=f(r_0)/f(r_1)>1$ for the concave parabola $f(x)=(x-r_3)(r_2-x)$ with vertex $-(r_0+r_1)/2<r_0$, while
$\rho_1>\rho_0$. This also gives $w_2=a(D_1^{-1}-D_0^{-1})/(P_1-P_0)>0$. Finally, cross-multiplying $w_2<w_1$, the
$a^2$-terms cancel and $D_0D_1D_3$ times the difference factors as
$2(r_0-r_1)(r_0-r_3)(r_1-r_3)(r_0r_3-r_1r_2)>0$.
\end{proof}
\begin{corollary}
(i) Outside the window, and for $m=0$, $\omega\neq0$ real, there is no real-frequency mode. (ii) $\omega=0$ lies outside the
window for $m\neq0$. (iii) At $w_2$ only $\operatorname{Im}e_t$ vanishes and at $w_1$ only $\operatorname{Im}e_1$; both
endpoints are then closed by the Wronskian and the injectivity of the transform (item 44).
\end{corollary}

\subsection{Inside the window: transform, dual potential, event lemma}
Inside the window the Riemann--Liouville transform $I_1^{\mu+1}$ maps a quasinormal mode injectively to a solution of a
dual Heun equation with the correct boundary branches (items 17, 55). The dual potential is exactly
$$L(z)=\frac Az+B-3z,\qquad A=t\big(X_0^2-(X_1+X_t)^2\big)>0,\qquad B=\lambda C_0+c_1,$$
with $C_0=3/(\Lambda(r_1-r_0)(r_2-r_3))$ and $c_1=1+t+3(1-\alpha)C_0$.

\begin{lemma}[the $c_1$ identity; item 89, \emph{cert}]
With $\Lambda=3$ (by scaling) and the exact relations $\sum\kappa_j=0$, $\sum r_j^2\kappa_j=0$,
$\sum r_j\kappa_j=-i\Xi\omega$ (residue identities), as rational functions
$n_1-(\lambda C_0+c_1)=-2(s^2-4)\frac{r_0^2+r_0r_1+r_0r_2-r_1r_2}{(r_0-r_1)(r_0+r_1+2r_2)}$ and $n_2+3=4-s^2$.
\end{lemma}
The angular bound $\lambda\ge3/2$ on the window (item 46) and $C_0>(2t-1)/(3(1-\alpha))$ (item 72, exact form
$C_0-\tfrac32\cdot\frac{2t-1}{3(1-\alpha)}=\frac{3(r_0+r_1)^2}{2\Lambda(r_1-r_0)(r_2-r_3)}$) give the floor
$B\ge5.5t-1.25$.

The two-current argument (items 8, 12, 17) reduces mode stability inside the window to the following.
\begin{definition}[event lemma]
For $t>1$, $z_0\in(1,t)$, $A>0$, $B\ge5.5t-1.25$, find $h$ on $(1,z_0)$ with $s=z_0-z$, $H=2\mu+s\mu'$ and
$$\text{(C)}\ \ s\,h\,\mu\le H,\qquad \text{(D)}\ \ \tfrac12\big(\mu'h+\mu(h'+h^2)\big)\le s(Lh+L').$$
\end{definition}
Regimes R1 ($p\le1$) and R2 ($p>1$, $2(z_0-1)^2(A/z_0+3)<t-1$), with $p=2(z_0-1)L(1)/(t-1)$, are closed by the explicit ODE
multiplier of item 47. In R3 we use the rational multiplier (item 104)
$$h=1.7\,g\,\frac{w}{1+w}+\frac{0.45\,(1-u)}{xz},\quad g=-\frac{L'}{L},\ w=8\kappa s^2,\ \kappa=\frac{L(z_0)}{\mu(z_0)},\ x=z-1,\ u=\frac x{z_0-1}.$$

\begin{proposition}[(C) in R3; item 105]
(C) holds for all $t,z_0,A$ and $B\ge5.5t-1.25$.
\end{proposition}
\begin{proof}
All three factors of $h$ are nonnegative, $w/(1+w)<1$, $1-u\le1$, so it suffices that
$1.7sg+0.45s/(xz)\le2+s/z+s/x-s/(t-z)$ (divide by $\mu$). Now $g-1/z=(6-B/z)/L$, so $g\le1/z$ when $B\ge6z$ and otherwise
the excess is at most $(0.5+1.25/t)/(2.5t-1.25)$; $0.7s/z+0.45s/(xz)-s/x\le-0.3s/z$; $s/(t-z)<1$; and where the excess is
nonzero $1.7s\cdot\text{excess}\le0.69$. Summing, the difference is at most $-0.3s/z+1+0.69-2<0$.
\end{proof}

\subsection{Regime R3: saturate (C), solve (D)}\label{sec:open}
Write $h=f'/f$. Then (D) $\iff \widetilde Df:=(\mu f')'-2s(Lf)'\le0$ and (C) $\iff h\le h_C:=H/(s\mu)=2/s+\mu'/\mu$.
\begin{proposition}[items 117--121]
For every R3 parameter point there is a $C^2$ function $f>0$ on $(1,z_0]$ with $f(1)=0$, (C) and (D) on $(1,z_0)$, and
$f'(z_0)>0$.
\end{proposition}
\begin{proof}
\emph{Construction.} On $(1,\hat z]$ let $f=F_C=\mu/s^2$, so (C) holds with equality; on $[\hat z,z_0]$ let $f$ solve
$\widetilde Df=0$ with matching data, so (D) holds with equality. Here $\hat z$ is the unique zero of
$d:=\widetilde DF_C/(2F_C)=N(S)/(2F_Cs^4)$, $S=z_0-z$.

\emph{(P1): $d$ has exactly one sign change.} $N$ is a sextic in $S$ with $N(0)=6\mu_0^2>0$ and vanishing $S^4$
coefficient. After $S=x_0v/(1+v)$ its seven coefficients $T_0,\dots,T_6$ satisfy $T_0,T_1>0$; $T_2,\dots,T_6$ are affine
in $(A,B)$ with negative coefficients; the feasible $(A,B)$-set is the quadrant $\{A\ge\max(0,A_{c2}),B\ge5.5t-1.25\}$;
and on it $T_4,T_5,T_6<0$ and $T_3\ge0\Rightarrow T_2>0$ (two-variable statements, exact Bernstein certificate).
Descartes' rule gives exactly one root.

\emph{(C) on the second piece.} With $p=\mu\exp(-\int2sL/\mu)$, the Wronskian $W=p(F_Cf'-F_C'f)$ satisfies
$W'=-(p/\mu)f\,\widetilde DF_C\le0$ and $W(\hat z)=0$, so $f'/f\le h_C$.

\emph{(P2): positivity and the seam.} At a zero of $f'$ with $f>0$, $\mu f''=2sL'f<0$, so $f'(z_0)>0$ implies $f'>0$ and
$f>0$ throughout. With $\Phi=\mu f'-2sLf$, $\Phi'=2Lf$; a first zero of $f'$ would force
$\Psi(z_1)\le0$ for a function $\Psi$ with $\Psi'<0$ and $\Psi(z_0)=\hat H/\hat s-2\int_{\hat z}^{z_0}s|L'|$, and
$\int_z^{z_0}s|L'|\le s^2(A/(2z^2)+3/2)$ reduces $\Psi(z_0)>0$ to $K(\hat z)>0$, $K=H-s^3|L'|$. Since $K$ and $N$ are
affine decreasing in $A$ (and $N$ in $B$), $K\le0\Rightarrow N(S;\max(A_K,0),5.5t-1.25)<0$, a three-variable statement
certified exactly; as $N\ge0$ on $[0,\hat S]$, $K(\hat z)>0$.
\end{proof}

\subsection{$\sigma=0$ and the upper half plane}
For $m=0$, $\sigma=0$: an energy identity with $V=-2\Lambda r^2-2(1-\alpha)-\lambda<0$ for $s=-2$, and for $s=+2$ the
Teukolsky--Starobinsky map, whose kernel is excluded by
$\det=[\lambda^2+6(1-\alpha)\lambda+8-4\alpha+8\alpha^2]^2>0$ (items 35, 43). Modes with $\operatorname{Im}\sigma>0$ are
excluded by continuity from Schwarzschild--de Sitter: QNMs depend continuously on $(\Lambda,a)$ (Petersen--Vasy, full
subextremal Fredholm theory) and none can cross the real axis by the above (item 36).

\section{Part B: ungauged linearised Einstein mode stability}
Throughout, $a\neq0$; for $a=0$ Theorem~B is Hintz--Vasy 2018. By Theorem~A, $\delta\Psi_0=\delta\Psi_4=0$ for every
generalised mode (item B.1). We use GHP with $\Th=D-p\epsilon-q\bar\epsilon$, $\eth=\delta-p\beta-q\bar\alpha$,
$\Th'=\Delta-p\gamma-q\bar\gamma$, $\eth'=\bar\delta-p\alpha-q\bar\beta$, $\rho'=-\mu$, $\tau'=-\pi$ (conventions fixed by
machine checks on KdS, item 92). On KdS, $\Psi_2=-M/(r-ia\cos\theta)^3$; put $\phi=r-ia\cos\theta\propto\Psi_2^{-1/3}$.

\subsection{Linearised Goldberg--Sachs}
\begin{lemma}[item 90]
Infinitesimal null rotations achieve $\delta\Psi_1=\delta\Psi_3=0$; then $\delta\kappa=\delta\sigma=\delta\kappa'=\delta\sigma'=0$.
\end{lemma}
\begin{proof}
To first order a null rotation about $n$ with parameter $d$ shifts $\delta\Psi_1$ by $3d\Psi_2$ and fixes
$\delta\Psi_0,\delta\Psi_2$; about $l$ similarly for $\delta\Psi_3$. As $\Psi_2\neq0$ the choice is algebraic. The
vacuum-$\Lambda$ Bianchi identities have no $\Lambda$ terms; linearising
$\Th\psi_1-\eth'\psi_0=-\tau'\psi_0+4\rho\psi_1-3\kappa\psi_2$ gives $\delta\kappa=0$, and
$\Th'\psi_0-\eth\psi_1=\rho'\psi_0-4\tau\psi_1+3\sigma\psi_2$ gives $\delta\sigma=0$; primes give the rest.
\end{proof}

\subsection{The Killing-spinor field and reality}
On a type D $\Lambda$-vacuum (with $\kappa=\sigma=\kappa'=\sigma'=0$) the GHP derivatives of the tetrad are as listed in
item 93(b) and the following hold (item 91; machine-checked on KdS and on the accelerating NUT Plebański--Demiański metric,
items 92, 97): $\Th\rho=\rho^2$, $\eth\rho=(\rho-\bar\rho)\tau$, $\Th\tau=(\tau-\bar\tau')\rho$, $\eth\tau=\tau^2$ and primes,
$\Th\rho'-\Th'\rho=\tau\bar\tau-\tau'\bar\tau'$, $\eth\tau'-\eth'\tau=\rho\bar\rho'-\bar\rho\rho'$, $\Th\tau'=\eth'\rho$,
$\Th'\tau=\eth\rho'$.
\begin{lemma}[items 93, 96]
With $A=\phi\rho$, $B=-\phi\rho'$, $C=\phi\tau$, $D=-\phi\tau'$, the field $\xi=Bl+An-Dm-C\bar m$ satisfies all ten
Killing equations as an algebraic consequence of these identities. On KdS, $\xi=-\Xi\,\partial_t$.
\end{lemma}
Hence for a perturbation in the gauge above, $\mathcal L_\xi h+\mathcal L_{\delta\xi}g=0$ with $\delta\xi$ a local linear
expression in $h$ commuting with $\mathcal L_{\partial_t},\mathcal L_{\partial_\varphi}$.
\begin{lemma}[reality for modes; item 96]
Let $J[h]=(\delta\xi[h]-\overline{\delta\xi[\bar h]})/(2i)$. Then $\mathcal L_{J[h]}g=0$, so $J[h]$ is a constant
combination of $\partial_t,\partial_\varphi$ (the local Killing algebra of KdS with $a\neq0$ is two-dimensional,
item 97 R2). For a mode with $(\sigma,m)\neq(0,0)$, $J[h]$ is a mode of the same type and so vanishes.
\end{lemma}
Consequently, to first order, $\phi\rho=\bar\phi\bar\rho$, $\phi\rho'=\bar\phi\bar\rho'$, $\phi\tau=-\bar\phi\bar\tau'$.

\subsection{The second Killing field}
Put $R=\operatorname{Re}\phi$, $I=\operatorname{Im}\phi$ (the Carter coordinates $r$, $a\cos\theta$). Reality gives
$\Th R=-A$, $\Th'R=B$, $\eth R=\eth'R=0$, $\Th I=\Th'I=0$, $\eth I=iC$, $\eth'I=-iD$.
\begin{lemma}[items 96, 97]
$\zeta=I^2(Bl+An)+R^2(Dm+C\bar m)$ satisfies all ten Killing equations by algebra (the $ln$, $m\bar m$, $lm$, $nm$
components are computed in item 96--97 and cancel by reality; the others are conjugates). On KdS
$\zeta=-\Xi(a^2\partial_t+a\partial_\varphi)$.
\end{lemma}

\subsection{Conclusion for $(\sigma,m)\neq(0,0)$}
From $\mathcal L_\xi h=-\mathcal L_{\delta\xi}g$ and $\mathcal L_\zeta h=-\mathcal L_{\delta\zeta}g$:
if $\sigma\neq0$, $h=\mathcal L_Xg$ with $X=\delta\xi/(i\sigma\Xi)$; if $\sigma=0$, $m\neq0$, $h=\mathcal L_Xg$ with
$X=-\delta\zeta/(im\Xi a)$. $X$ is smooth across the horizons because $\xi,\zeta,\delta\xi,\delta\zeta$ are boost-spin
invariant natural constructions from the Weyl spinor (item 97 R3). HPV's modes with zero $T$-frequency and $m\neq0$
have $\partial_t$-frequency $\Omega m\neq0$.

\subsection{$\sigma=m=0$}
\begin{proposition}[items 98, 102, 103]
A stationary axisymmetric $h$ with $\delta\Psi_0=\delta\Psi_4=0$ is $c_M\partial_Mg+c_a\partial_ag+\mathcal L_Xg$.
\end{proposition}
\begin{proof}[Sketch]
The constants in $J[h]$ are removed by a phase of $\phi$ and by the acceleration tangent, which is the only
Plebański--Demiański tangent making $\xi$ complex (item 95, exact nonlinear check). Then $\xi$, $\zeta$ are first-order
Killing fields, they commute, $\nabla R\perp\nabla I$, $\Delta_r:=-2AB|\phi|^2$ depends on $R$ only and
$\Delta_x:=-2CD|\phi|^2$ on $I$ only, and in adapted coordinates the perturbed metric is Carter's ansatz to first order.
On that ansatz the $\Lambda$-vacuum equations are linear in $(\Delta_r,\Delta_x)$, with solutions exactly
$\Delta_r=-\tfrac\Lambda3R^4+\varepsilon R^2-2mR+\gamma$, $\Delta_x=-\tfrac\Lambda3I^4-\varepsilon I^2+2nI+\gamma$
(item 103, \emph{cert}). Modulo the scaling, the variations are $\delta M,\delta a,\delta n$; with acceleration, four.
Axis regularity (determinant $-64mp_0^4/c\neq0$) removes NUT and acceleration.
\end{proof}

\subsection{Generalised modes}
For $(\sigma,m)\neq(0,0)$, induction on $k$: the top coefficient is a mode, hence $\mathcal L_{X_1}g$, and
$Y=t_*^{k-1}X_1/(k-1)!$ lowers the order. For $\sigma=m=0$, linear growth of $(\delta M,\delta a)$ is excluded by the
Iyer--Wald charges: they are surface-independent, vanish on every $\mathcal L_Xg$ by covariance of the Noether charge
(item 101), and $\partial(E,J)/\partial(M,a)$ has determinant $M/\Xi^4\neq0$ for $E=M/\Xi^2$, $J=aM/\Xi^2$ (item 100).

\section{What remains}
\begin{enumerate}
\item An independent check of the whole argument, in particular of Part~B (the in-house hostile passes are items 97
and 102) and of the R3 construction (items 117--121).
\end{enumerate}

\end{document}
